% ======================================================================
% 1. TARGET CONFERENCE SELECTION
% ======================================================================
\newif\ifIEEE
\newif\ifACM
\newif\ifUSENIX

% ----------------------------------------------------------------------
\ifdefined\ForceACM
    \ACMtrue
\else\ifdefined\ForceUSENIX
    \USENIXtrue
\else\ifdefined\ForceIEEE
    \IEEEtrue
\else
    \ACMtrue % Default to ACM if no specific conference is forced
\fi\fi\fi
% ----------------------------------------------------------------------

% ======================================================================
% 2. DOCUMENT CLASS & SPECIFIC PACKAGES
% ======================================================================
\ifIEEE
    \documentclass[conference, 10pt]{IEEEtran}
    \IfFileExists{cmap.sty}{\usepackage{cmap}}{}
    \IEEEoverridecommandlockouts
    \usepackage{cite}
    \usepackage{amsmath,amssymb,amsfonts}
    \usepackage{algorithmic}
\fi

\ifACM
    \documentclass[sigconf,10pt]{acmart}
    \renewcommand\footnotetextcopyrightpermission[1]{} % remove copyright stuff
    \settopmatter{printfolios=true,printacmref=false}  % page num
\fi

\ifUSENIX
    \documentclass[letterpaper,twocolumn,10pt]{article}
    \IfFileExists{cmap.sty}{\usepackage{cmap}}{}
    \usepackage{usenix-2020-09}
    \usepackage{amsmath,amssymb}
\fi

% Improve copy/paste/searchability in pdfTeX builds. ACM already enables these
% hooks in acmart.cls, so only add them for the other targets here.
\ifACM\else
    \ifdefined\pdfgentounicode
        \IfFileExists{glyphtounicode.tex}{\input{glyphtounicode}}{}
        \pdfgentounicode=1
    \fi
\fi

% ======================================================================
% 3. COMMON PACKAGES
% ======================================================================
\usepackage{graphicx}
\usepackage{caption}
\usepackage{subcaption}
\graphicspath{{./figures/}}
\usepackage{tikz}
\usetikzlibrary{arrows.meta,calc,positioning,decorations.pathreplacing,patterns,fit,backgrounds}
\usepackage{xcolor}
\usepackage{tcolorbox}
\usepackage[ruled,lined,noend]{algorithm2e}
\usepackage{listings}
\lstdefinestyle{pseudocode}{
    language=Python,
    basicstyle=\small\ttfamily,
    keywordstyle=\bfseries,
    commentstyle=\itshape\color{black!55},
    showstringspaces=false,
    columns=fullflexible,
    keepspaces=true,
    xleftmargin=1em,
    numbers=none,
    escapeinside={(@}{@)},
    morekeywords={send,await,to,None,True,False,match},
    deletekeywords={type},
    aboveskip=0pt,
    belowskip=0pt,
    literate=
        {!=}{{$\neq$}}1
        {<=}{{$\leq$}}1
        {>=}{{$\geq$}}1,
}
\usepackage{booktabs}
\usepackage{pgfplots}
\pgfplotsset{compat=1.18}
\usepgfplotslibrary{statistics}

% Preserve most of microtype's typography benefits while avoiding the spacing
% adjustments most likely to confuse PDF text extraction.
\makeatletter
\@ifpackageloaded{microtype}{\microtypesetup{tracking=false,spacing=false}}{}
\makeatother

% Prevent overfull \hbox in narrow two-column layout
\emergencystretch=1em
% Mildly reduce hyphenation without introducing loose spacing
\tolerance=1000
\hyphenpenalty=200
\hbadness=1000

\begin{document}

\newcommand{\red}[1]{\textcolor{red}{#1}}
\newcommand{\green}[1]{\textcolor{green}{#1}}
\newcommand{\blue}[1]{\textcolor{blue}{#1}}

\newcommand{\tlide}{$\sim$}
\newcommand{\mul}{$\times$}
\newcommand{\eg}{\emph{e.g.}}
\newcommand{\etal}{\emph{et al.}}
\newcommand{\summary}[1]{\noindent\textbf{#1}}

\newcommand*\circled[1]{\tikz[baseline=(char.base)]{
            \node[shape=circle,fill,inner sep=1pt] (char) {\textcolor{white}{#1}};}}

\newenvironment{tightitemize}%
 {\begin{list}{$\bullet$}{%
 		\setlength{\leftmargin}{10pt}
        \setlength{\itemsep}{0pt}%
        \setlength{\parsep}{0pt}%
        \setlength{\topsep}{0pt}%
        \setlength{\parskip}{0pt}%
        }%
 }%
{\end{list}}

% ------------------------------------------------------------------
% Raft-style protocol summary figures  (tcolorbox-based)
% ------------------------------------------------------------------
\definecolor{prothdrcolor}{HTML}{1F4799}
\definecolor{protsubcolor}{HTML}{A3001F}

% Helvetica for protocol-figure headers (matches the Raft cheatsheet)
\newcommand{\protsfamily}{\fontfamily{phv}\selectfont}

% Layout dimensions
\newlength{\protboxgap}              % horizontal gap between columns
\setlength{\protboxgap}{8pt}
\newlength{\protboxvgap}             % vertical gap between boxes
\setlength{\protboxvgap}{8pt}
\newlength{\protcolwidth}            % column width (set by \protsetup)
\newlength{\protfieldwd}             % field-name column width (set by \protsetup)

% Compute column and field widths.  Call once per figure,
% inside the \footnotesize group so em is relative to body text.
% The -1pt avoids dimexpr rounding that can exceed \textwidth.
\newcommand{\protsetup}{%
  \setlength{\protcolwidth}{\dimexpr(\textwidth - \protboxgap)/2 - 1pt\relax}%
  \setlength{\protfieldwd}{6.5em}%
}

% Protocol box: rounded corners, colored border = header color.
% Width defaults to \linewidth, which is set by the enclosing minipage.
\newtcolorbox{protbox}[1]{%
  title={#1},
  colframe=prothdrcolor,
  colback=white,
  colbacktitle=prothdrcolor,
  coltitle=white,
  fonttitle=\protsfamily\bfseries,
  halign title=center,
  rounded corners,
  arc=2pt,
  boxrule=0.8pt,
  left=4pt,
  right=4pt,
  top=2pt,
  bottom=3pt,
  toptitle=1pt,
  bottomtitle=1pt,
  titlerule=0pt,
  nobeforeafter,
  box align=top,
  after=\par,
  before upper={\setlength{\parskip}{0pt}\setlength{\parindent}{0pt}},
}

% Sub-header in dark red, Helvetica bold
\newcommand{\protsubhdr}[1]{%
  \par\smallskip\noindent\textcolor{protsubcolor}{\protsfamily\bfseries #1}\par\noindent}

% Two-column field table — fixed left column for names,
% flexible right column for descriptions.  Keeps descriptions
% aligned across sub-headers and across boxes.
\newenvironment{protfields}{%
  \begin{tabular}{@{}p{\protfieldwd}@{\hspace{0.5em}}p{\dimexpr\linewidth-\protfieldwd-0.5em}@{}}%
}{%
  \end{tabular}%
}

% Tight numbered list for receiver-implementation steps
\newcounter{rpcstep}
\newenvironment{protimpl}{%
  \setcounter{rpcstep}{0}%
  \begin{list}{\arabic{rpcstep}.}{%
    \usecounter{rpcstep}%
    \setlength{\leftmargin}{1.5em}%
    \setlength{\rightmargin}{0pt}%
    \setlength{\itemsep}{1pt}%
    \setlength{\parsep}{0pt}%
    \setlength{\topsep}{2pt}%
    \setlength{\partopsep}{0pt}%
  }%
}{%
  \end{list}%
}

% Tight bulleted list for proposer rules
\newenvironment{protrules}{%
  \begin{list}{$\bullet$}{%
    \setlength{\leftmargin}{1.5em}%
    \setlength{\rightmargin}{0pt}%
    \setlength{\itemsep}{1pt}%
    \setlength{\parsep}{0pt}%
    \setlength{\topsep}{2pt}%
    \setlength{\partopsep}{0pt}%
  }%
}{%
  \end{list}%
}

% ------------------------------------------------------------------
% Compact space-time figures (space on y-axis, time on x-axis)
% ------------------------------------------------------------------
\definecolor{streqfill}{HTML}{FFF2CC}
\definecolor{strspfill}{HTML}{DDEBF7}
\definecolor{stcommitfill}{HTML}{D9EAD3}
\definecolor{stmodefill}{HTML}{E8F1FF}
\definecolor{stnotefill}{HTML}{F3F3F3}

\tikzset{
  stproc/.style={font=\footnotesize\bfseries, anchor=east},
  stlane/.style={densely dotted, gray!55, thin},
  streq/.style={->, line width=0.85pt, shorten <=1pt, shorten >=1pt},
  strsp/.style={->, dashed, line width=0.85pt, shorten <=1pt, shorten >=1pt},
  staxis/.style={->, line width=0.85pt},
  streqbox/.style={
    draw=black!50,
    fill=streqfill,
    rounded corners=1.5pt,
    inner xsep=2pt,
    inner ysep=1.2pt,
    font=\scriptsize\protsfamily,
    align=center,
  },
  strspbox/.style={
    draw=black!50,
    fill=strspfill,
    rounded corners=1.5pt,
    inner xsep=2pt,
    inner ysep=1.2pt,
    font=\scriptsize\protsfamily,
    align=center,
  },
  stcommitbox/.style={
    draw=green!50!black,
    fill=stcommitfill,
    rounded corners=1.5pt,
    inner xsep=2pt,
    inner ysep=1.2pt,
    font=\scriptsize\protsfamily,
    align=center,
  },
  stnotebox/.style={
    draw=gray!40,
    fill=stnotefill,
    rounded corners=1.5pt,
    inner xsep=2pt,
    inner ysep=1.2pt,
    font=\scriptsize\protsfamily,
    align=center,
  },
  stbluebox/.style={
    draw=blue!35,
    fill=blue!6,
    rounded corners=1.5pt,
    inner xsep=2.5pt,
    inner ysep=1.5pt,
    font=\scriptsize,
    text=blue!70!black,
    align=center,
  },
  stwarnbox/.style={
    draw=red!45!black,
    fill=red!8,
    rounded corners=1.5pt,
    inner xsep=2pt,
    inner ysep=1.2pt,
    font=\scriptsize,
    text=red!70!black,
    align=center,
  },
  stmodebox/.style={
    draw=blue!35,
    fill=stmodefill,
    rounded corners=1.5pt,
    inner xsep=2pt,
    inner ysep=1.2pt,
    font=\scriptsize\itshape,
    text=blue!65!black,
    align=center,
  },
  stsection/.style={font=\scriptsize\itshape, text=gray!75!black, align=center},
  stquiet/.style={font=\scriptsize, text=gray!70!black},
  stfail/.style={font=\scriptsize, text=red!65!black},
}



% ======================================================================
% 4. TITLE & AUTHORS
% ======================================================================
\title{Fast CASPaxos: Leaderless, Linearizable, Logless}

\ifIEEE
    \author{
        \IEEEauthorblockN{Reuben Bond}
        \IEEEauthorblockA{
            reuben.bond@microsoft.com
        }
    }
\fi

\ifACM
    \author{Reuben Bond}
    \email{reuben.bond@microsoft.com}
\fi

\ifUSENIX
    \date{} % don't want date printed
    \author{
      {\rm Reuben Bond}\\
      reuben.bond@microsoft.com
    }
\fi

% ======================================================================
% 5. ABSTRACT & MAKETITLE ORDERING
% ======================================================================
\ifACM
    % ACM requires abstract BEFORE maketitle
    \begin{abstract}
Paxos and its variants underpin the replicated state machines at the heart of many distributed systems.
CASPaxos simplifies single-decree Paxos into a replicated, rewritable, linearizable register that requires only two metadata fields per key.
Like most consensus protocols, CASPaxos is \emph{leadered}: a proposer must first obtain exclusive rights before committing a value, giving a fast path of one round-trip (1~RTT) only when leadership is uncontended.
Fast Paxos introduced \emph{leaderless} fast rounds in which any proposer may commit in 1~RTT without prior leader election, at the cost of larger quorums and a recovery path when concurrent proposals conflict.

In this paper we present \emph{Fast CASPaxos}, a protocol that blends the Fast Paxos fast-round optimization with CASPaxos.
By letting each piggybacked prepare choose between a shared fast ballot and a proposer-owned classic ballot, Fast CASPaxos allows the same implementation to switch between leadered and leaderless operation on a per-operation basis.
We describe the protocol, argue its correctness, present a TLA+ specification verified by model checking, and evaluate its performance through deterministic simulation across LAN and geo-distributed topologies with fault injection and linearizability checking.
\end{abstract}

% ---------------------------------------------
% Keywords for different formats
% ---------------------------------------------
% 1. IEEE: IEEEkeywords
\ifdefined\IEEEkeywords
    \begin{IEEEkeywords}
    distributed consensus, Paxos, leaderless consensus, replicated state machines
    \end{IEEEkeywords}
\fi

% 2. ACM: ccsdesc
\ifdefined\ccsdesc
    \ccsdesc[500]{Theory of computation~Distributed algorithms}
    \keywords{distributed consensus, Paxos, leaderless consensus, replicated state machines}
\fi

% 3. USENIX
\ifdefined\USENIXtrue
\ifUSENIX
    \vspace{0.5cm}
    \noindent \textbf{Keywords:} distributed consensus, Paxos, leaderless consensus, replicated state machines
\fi
\fi

    \maketitle
    \pagenumbering{arabic}
    \pagestyle{plain} % to remove upper notes
\else
    % IEEE and USENIX require maketitle BEFORE abstract
    \maketitle
    \begin{abstract}
Paxos and its variants underpin the replicated state machines at the heart of many distributed systems.
CASPaxos simplifies single-decree Paxos into a replicated, rewritable, linearizable register that requires only two metadata fields per key.
Like most consensus protocols, CASPaxos is \emph{leadered}: a proposer must first obtain exclusive rights before committing a value, giving a fast path of one round-trip (1~RTT) only when leadership is uncontended.
Fast Paxos introduced \emph{leaderless} fast rounds in which any proposer may commit in 1~RTT without prior leader election, at the cost of larger quorums and a recovery path when concurrent proposals conflict.

In this paper we present \emph{Fast CASPaxos}, a protocol that blends the Fast Paxos fast-round optimization with CASPaxos.
By letting each piggybacked prepare choose between a shared fast ballot and a proposer-owned classic ballot, Fast CASPaxos allows the same implementation to switch between leadered and leaderless operation on a per-operation basis.
We describe the protocol, argue its correctness, present a TLA+ specification verified by model checking, and evaluate its performance through deterministic simulation across LAN and geo-distributed topologies with fault injection and linearizability checking.
\end{abstract}

% ---------------------------------------------
% Keywords for different formats
% ---------------------------------------------
% 1. IEEE: IEEEkeywords
\ifdefined\IEEEkeywords
    \begin{IEEEkeywords}
    distributed consensus, Paxos, leaderless consensus, replicated state machines
    \end{IEEEkeywords}
\fi

% 2. ACM: ccsdesc
\ifdefined\ccsdesc
    \ccsdesc[500]{Theory of computation~Distributed algorithms}
    \keywords{distributed consensus, Paxos, leaderless consensus, replicated state machines}
\fi

% 3. USENIX
\ifdefined\USENIXtrue
\ifUSENIX
    \vspace{0.5cm}
    \noindent \textbf{Keywords:} distributed consensus, Paxos, leaderless consensus, replicated state machines
\fi
\fi

\fi

% ======================================================================
% 6. MAIN BODY
% ======================================================================
% ==============================================================================
\section{Introduction}
\label{sec:introduction}
% ==============================================================================

Distributed consensus allows a set of processes to agree on a single value despite failures.
Paxos~\cite{lamport_tcs98,lamport_sigact01} and its variants provide reliable solutions to this problem and often underpin the replicated state machines at the heart of industrial distributed systems~\cite{schneider_cs90}.

Classic Paxos operates in two phases---\emph{prepare} and \emph{accept}---each requiring a quorum of acceptors.
Multi-Paxos amortizes the cost of the prepare phase by electing a stable leader, allowing committed values in a single round-trip as long as the leader is uncontended.
However, the leader is also a single point of contention and a potential performance bottleneck.
If the leader fails, the system must detect the failure and elect a new leader before progress can resume.

Fast Paxos~\cite{lamport_msr05} introduced the concept of \emph{fast rounds}, in which any proposer may send a value directly to acceptors without first acquiring leadership.
If sufficiently many acceptors receive the same value, it commits in one round-trip.
When concurrent proposals conflict, the protocol falls back to a coordinated recovery round.
The price is a larger quorum requirement: fast quorums must be large enough that any later classic recovery quorum can identify which fast-round value, if any, was committed.

CASPaxos~\cite{rystsov_caspaxos} simplifies single-decree Paxos into a replicated, rewritable, linearizable register with compare-and-swap semantics.
It requires only two metadata fields per register and supports a piggybacked-prepare optimization that amortizes leader election across successive operations.

In this paper we present \emph{Fast CASPaxos}, which applies the Fast Paxos fast-round optimization to CASPaxos.
The key insight is that the piggybacked ballot in each accept message can be either a shared fast ballot or a proposer-owned classic ballot, allowing the protocol to switch between leadered and leaderless operation on a per-operation basis.
This gives deployments the flexibility to optimize for their current workload: leaderless mode for geo-distributed, low-conflict scenarios where 1~RTT commits from any proposer are attractive, and leadered mode when contention is high and a stable leader provides better throughput.

We make the following contributions:
\begin{itemize}
    \item We describe the Fast CASPaxos protocol and its adaptive mode-switching mechanism (Section~\ref{sec:fast-caspaxos}).
    \item We argue correctness and present a TLA+ specification verified by model checking (Section~\ref{sec:correctness}, Appendix~\ref{sec:tla-specification}).
    \item We evaluate performance through deterministic simulation with fault injection and linearizability checking across LAN and geo-distributed topologies (Section~\ref{sec:evaluation}).
\end{itemize}

% ==============================================================================
\section{Background}
\label{sec:background}
% ==============================================================================

% ------------------------------------------------------------------------------
\subsection{Paxos and Fast Paxos}
\label{sec:paxos-fast-paxos}
% ------------------------------------------------------------------------------

Paxos~\cite{lamport_tcs98,lamport_sigact01} solves single-decree consensus among a set of acceptors coordinated by one or more proposers.
The protocol proceeds in two phases.
In the \emph{prepare} phase, a proposer selects a ballot number and sends it to a quorum of acceptors; each acceptor promises not to accept proposals with lower ballot numbers and returns any value it has already accepted.
In the \emph{accept} phase, the proposer sends a value to a quorum of acceptors using the ballot from the prepare phase; if the quorum accepts, the value is committed.
Any two quorums must intersect, which is satisfied by majority quorums of size $\lfloor N/2 \rfloor + 1$ for $N$ acceptors.

Multi-Paxos optimizes for sequences of decisions by electing a stable leader.
Once a leader has completed a prepare phase, it can commit subsequent values in a single round-trip by reusing its ballot.
This amortized cost makes Multi-Paxos attractive for high-throughput workloads, but the leader becomes a bottleneck and a single point of failure.

Fast Paxos~\cite{lamport_msr05} introduced \emph{fast rounds} in which any proposer may send a value directly to acceptors without prior leader election.
In a fast round, multiple proposers may propose different values under the same ballot.
If a sufficiently large \emph{fast quorum} of acceptors all accept the same value, it commits in one round-trip.
When acceptors receive conflicting values, the protocol falls back to a coordinated classic recovery round.

The fast quorum must be a supermajority---at least $\lceil 3N/4 \rceil$ acceptors---to ensure that any later classic recovery quorum can determine which value, if any, was committed during the fast round.
This is because a classic quorum that reads the state of a fast round must see enough copies of the committed value to distinguish it from any competing value.

Flexible Paxos~\cite{howard_opodis16} showed that Paxos's quorum intersection requirement can be relaxed: intersection is only needed between phases, not within them.
Howard et al.~\cite{howard_icdcn21} extended this result to Fast Paxos, showing that the quorum intersection requirements of fast rounds can also be weakened.

% ------------------------------------------------------------------------------
\subsection{CASPaxos}
\label{sec:caspaxos}
% ------------------------------------------------------------------------------

CASPaxos~\cite{rystsov_caspaxos} simplifies single-decree Paxos into a protocol for a replicated, rewritable, linearizable register.
Rather than choosing a single immutable value, CASPaxos allows proposers to apply an \emph{update function} to the register's current value, providing compare-and-swap (CAS) semantics.

Each acceptor maintains only two fields per register: a \emph{promised ballot} (the highest ballot it has promised to honor) and an \emph{accepted ballot/value} pair (the last value it accepted and the ballot under which it was accepted).
This minimal metadata footprint---typically 16 to 32 bytes---makes CASPaxos attractive for systems that manage many fine-grained replicated keys.

CASPaxos supports a \emph{piggybacked-prepare} optimization: each accept message includes the ballot for the \emph{next} prepare phase.
If the acceptor processes the accept successfully, it also promises the piggybacked ballot, effectively amortizing the prepare cost.
This allows a proposer that ``owns'' the piggybacked ballot to commit its next value in a single round-trip, much like Multi-Paxos's leader optimization but without a separate leader-election mechanism.

% ==============================================================================
\section{Fast CASPaxos}
\label{sec:fast-caspaxos}
% ==============================================================================

Fast CASPaxos extends CASPaxos with shared fast ballots.
The piggybacked prepare becomes an atomic promise for the next ballot, which may be proposer-owned or shared.

\begin{figure*}[t]
\centering
{\footnotesize
\protsetup
\noindent
\begin{minipage}[t]{\protcolwidth}
\begin{protbox}{State}
\protsubhdr{Proposer state \textit{(volatile):}}
\begin{protfields}
\textbf{id} & unique proposer identifier ($\textit{id} > 0$)\\
\textbf{round} & monotonically increasing, initially 0\\
\textbf{prepared} & most recently prepared ballot, initially $(1, 0)$\\
\textbf{cachedVal} & last decided value, initially \texttt{null}\\
\end{protfields}
\protsubhdr{Acceptor state \textit{(persistent):}}
\begin{protfields}
\textbf{promised} & highest ballot promised, initially $(1, 0)$\\
\textbf{accepted} & last accepted ballot, or \texttt{null}\\
\textbf{value} & last accepted value, or \texttt{null}\\
\end{protfields}
\end{protbox}
\vspace{\protboxvgap}
\begin{protbox}{Rules for Proposers}
\protsubhdr{Quorums:}
$q_p = q_c = \lfloor N/2 \rfloor + 1$ (prepare and classic accept),\; $q_f = \lceil 3N/4 \rceil$ (fast accept).
\protsubhdr{Procedure:}
\textit{To execute an operation with update function $f$:}
\begin{protimpl}
\item \textbf{Read.}\; Obtain ballot $b$ and \textit{cachedVal}:
  \begin{protrules}
  \item If \textit{prepared} is not \texttt{null}: let $b = \textit{prepared}$;\; set \textit{prepared} to \texttt{null}.
  \item Else: increment \textit{round}; let $b = (\textit{round}, \textit{id})$; send \textsc{Prepare}($b$) to all acceptors. Wait for $q_p$~promises; set \textit{cachedVal} per Value Recovery below. If fewer than $q_p$ promises: set \textit{round} above highest round from any rejections; retry step 1.
  \end{protrules}
\item \textbf{Modify.}\; let $\textit{newVal} = f(\textit{cachedVal})$.
\item \textbf{Write.}\; Choose \textit{nextBallot}: $(b.r{+}1,\, 0)$ to prepare a fast round;\; $(b.r{+}1,\, \textit{id})$ to prepare a classic round.
\item Send \textsc{Accept}($b$, \textit{newVal}, \textit{nextBallot}) to all acceptors. Wait for $q_f$~responses if $b$ is a fast ballot; $q_c$ if $b$ is a classic ballot.
\item On \textsc{Accept} response:
  \begin{protrules}
  \item Quorum success: $\textit{cachedVal} = \textit{newVal}$;\; if a quorum of successful responses reports $\textit{promised} \ge \textit{nextBallot}$ then set $\textit{prepared} = \textit{nextBallot}$, else set \textit{prepared} to \texttt{null}.
  \item Otherwise (rejection or timeout): set \textit{prepared} to \texttt{null}; set \textit{round} above highest round from any rejections; retry from step 1.
  \end{protrules}
\end{protimpl}
\protsubhdr{Value Recovery:}
\textit{When \textsc{Prepare} returns accepted values from a prepare quorum:}
\begin{protrules}
\item If no response contains an \textit{accepted} ballot: set \textit{cachedVal} to \texttt{null}.
\item If the highest \textit{accepted} ballot is \textbf{classic}: $\textit{cachedVal}$ = the corresponding \textit{value}. Only one value can exist at a classic ballot.
\item If the highest \textit{accepted} ballot is \textbf{fast}: tally responses by \textit{value} at that ballot. Let \texttt{maxCount} be the highest tally and set \textit{cachedVal} to the unique maximum, or any tied maximum.
\end{protrules}
\end{protbox}
\end{minipage}%
\hfill
\begin{minipage}[t]{\protcolwidth}
\begin{protbox}{Ballots}
A ballot is a tuple $(r, p)$ where $r$ is a round number and $p$ is the proposer's \textbf{id}. \textbf{Classic} ($p \neq 0$): only proposer $p$ may use; a classic ballot corresponds to at most one value. \textbf{Fast} ($p = 0$): any proposer may propose.
\par\smallskip\noindent
Ordering: $(r_1, p_1) > (r_2, p_2)$ iff $r_1 > r_2$, or $r_1{=}r_2$ and $p_1 > p_2$.
\end{protbox}
\vspace{\protboxvgap}
\begin{protbox}{Prepare RPC}
\textit{Invoked by proposer to read the current value and acquire exclusive write access.}
\protsubhdr{Arguments:}
\begin{protfields}
\textbf{ballot} & ballot $(r, p)$ for this operation\\
\end{protfields}
\protsubhdr{Results:}
\begin{protfields}
\textbf{success} & \texttt{Promise(accepted, value)}\\
\textbf{failure} & \texttt{Reject(promised)}\\
\end{protfields}
\protsubhdr{Acceptor implementation:}
\begin{protimpl}
\item If $\textit{ballot} < \textit{promised}$, reply \texttt{Reject(promised)}.
\item $\textit{promised} = \textit{ballot}$.
\item Reply \texttt{Promise(accepted, value)}.
\end{protimpl}
\end{protbox}
\vspace{\protboxvgap}
\begin{protbox}{Accept RPC}
\textit{Invoked by proposer to commit a value to the register.}
\protsubhdr{Arguments:}
\begin{protfields}
\textbf{ballot} & ballot to write under\\
\textbf{value} & the value to write\\
\textbf{nextBallot} & optional piggybacked prepare; \texttt{null} skips it\\
\end{protfields}
\protsubhdr{Results:}
\begin{protfields}
\textbf{success} & \texttt{Accepted(promised)}\\
\textbf{failure} & \texttt{Reject(promised)}\\
\end{protfields}
\protsubhdr{Acceptor implementation:}
\begin{protimpl}
\item If $\textit{ballot} < \textit{promised}$, reply \texttt{Reject(promised)}.
\item If $\textit{ballot} = \textit{accepted}$:
  \begin{protrules}
  \item If $\textit{value} = v$: reply \texttt{Accepted(promised)} (state unchanged).
  \item If $\textit{value} \neq v$: reply \texttt{Reject(promised)}.
  \end{protrules}
\item $\textit{accepted} = \textit{ballot}$;\; $\textit{value} = v$.
\item $\textit{promised} = \max(\textit{ballot},\, \textit{nextBallot})$.
\item Reply \texttt{Accepted(promised)}.
\end{protimpl}
\end{protbox}
\end{minipage}%
}% footnotesize
\caption{Fast CASPaxos summary. Relative to CASPaxos (Figure~\ref{fig:caspaxos-summary}), ballots with $p = 0$ are shared fast rounds, fast accepts require $\lceil 3N/4 \rceil$ acceptors, fast ballots use first-write-wins with idempotent duplicates, and Value Recovery chooses the unique maximum or any tied maximum.}
\label{fig:fast-caspaxos-summary}
\end{figure*}


% ------------------------------------------------------------------------------
\subsection{Piggybacked Prepare}
\label{sec:piggybacked-prepare}
% ------------------------------------------------------------------------------

In CASPaxos, a successful \textsc{Accept} may carry a piggybacked prepare for the next ballot.
Fast CASPaxos keeps this mechanism but lets the next ballot be either
\begin{itemize}
    \item $(r+1, p)$: a proposer-owned classic ballot;
    \item $(r+1, 0)$: a shared fast ballot.
\end{itemize}

The piggybacked prepare affects only later rounds: it adds no state and does not change the current accept.

\paragraph{Register initialization.}
Acceptors are initialized with promised ballot $(1, 0)$, so the first operation may attempt a fast round immediately.
If it collides, recovery proceeds through an ordinary classic round.

\paragraph{Partial piggybacked-prepare delivery.}
A proposer may reuse a piggybacked ballot only if a quorum of successful \textsc{Accept} responses confirms that promise.
Otherwise it discards the candidate next ballot and falls back to \textsc{Prepare} on its next operation.

% ------------------------------------------------------------------------------
\subsection{Fast Ballots}
\label{sec:fast-ballots}
% ------------------------------------------------------------------------------

Multiple proposers may concurrently submit different values under the same fast ballot, and different acceptors may accept different values.
Safety requires that at most one value be \emph{chosen}, meaning accepted by a fast quorum.
Fast CASPaxos enforces this by having an acceptor accept the first valid request it receives for a fast ballot and reject later requests for that ballot with a different value.
If concurrent proposals split the acceptors, several values may be accepted, but none is chosen unless one reaches a fast quorum.
Because fast quorums intersect, two different values cannot both be chosen in the same fast ballot; if none reaches a fast quorum, a later classic round recovers a safe value from the highest accepted ballot it sees.

\subsubsection{Idempotent Write for Identical Values}
When multiple proposers submit \emph{identical} values under the same fast ballot---as in almost-everywhere-agreement workloads such as Rapid cluster membership~\cite{suresh_atc18}---strict first-write-wins would cause unnecessary rejections.
We therefore let an acceptor acknowledge an \textsc{Accept} request that exactly matches its current accepted ballot and value.
This does not change state and is safe by Section~\ref{sec:correctness}.

\subsubsection{Value Recovery}
\label{sec:value-recovery}
When a proposer runs \textsc{Prepare}, it must recover a value from the highest accepted ballot reported by the prepare quorum.
The recovery rule is best phrased in terms of the values that \emph{could have been decided} at that highest ballot:

\begin{itemize}
    \item \textbf{No prior value.} If no acceptor in the quorum has accepted any value, the proposer is free to propose any value.
    \item \textbf{Classic ballot.} If the highest accepted ballot is a classic ballot, its value is the only possible decided value at that ballot.
    \item \textbf{Fast ballot.} If the highest accepted ballot is a fast ballot, several values may appear at that ballot. In the fixed-quorum setting, the possible decided values are exactly those tied for maximum support in the prepare responses at that ballot. A unique maximum must be recovered; any tied maximum is safe.
\end{itemize}

\noindent A value chosen in a fast ballot was accepted by $q_f$ acceptors; any later prepare quorum intersects that fast quorum in at least $q_f + q_p - N$ acceptors, so any chosen value appears in the recovery evidence.\footnote{This is the counting-quorum specialisation of the \emph{possibleValues} function from the epochs-by-recovery framework~\cite{howard_tr19}.}
A uniquely chosen value therefore appears strictly more often than any competitor, while an exact tie leaves the prepare quorum unable to distinguish among the tied values.

\subsubsection{Conflict Recovery Walkthrough}
Consider $N = 4$ acceptors with fast quorum $q_f = \lceil 3 \cdot 4 / 4 \rceil = 3$, all prepared for fast ballot $(1, 0)$.
Suppose $P_1$ proposes $v_a$ and $P_2$ proposes $v_b$, and the acceptors split 2/2: $a_1,a_2$ accept $v_a$, $a_3,a_4$ accept $v_b$, and each acceptor rejects the other value under the first-write-wins rule.
No value reaches a fast quorum, so no value is chosen.
$P_1$ then runs \textsc{Prepare}($(2, 1)$) and collects responses from, say, $\{a_1, a_2, a_3\}$, which report $v_a$ twice and $v_b$ once at ballot $(1, 0)$.
Because $v_a$ is the unique maximum at the highest fast ballot seen, $P_1$ must recover $v_a$, send \textsc{Accept}($(2, 1),\, v_a$), and commit after a classic quorum of $q_c = 3$ confirms.
Recovery therefore adds two message delays beyond the failed fast round, for a total of 3~RTT from the client's perspective.
Figure~\ref{fig:collision-recovery} illustrates this scenario.

\begin{figure}[t]
\centering
\resizebox{\columnwidth}{!}{%
\begin{tikzpicture}[x=1.12cm, y=1.20cm, font=\footnotesize, >=Stealth]
\def\xEnd{8.8}
\def\yPa{3.00}
\def\yPb{2.40}
\def\yA{1.80}
\def\yB{1.20}
\def\yC{0.60}
\def\yD{0.0}

\node[stproc] at (0,\yPa) {$P_1$};
\node[stproc] at (0,\yPb) {$P_2$};
\node[stproc] at (0,\yA) {$a_1$};
\node[stproc] at (0,\yB) {$a_2$};
\node[stproc] at (0,\yC) {$a_3$};
\node[stproc] at (0,\yD) {$a_4$};

\draw[stlane] (0.25,\yPa) -- (\xEnd,\yPa);
\draw[stlane] (0.25,\yPb) -- (\xEnd,\yPb);
\draw[stlane] (0.25,\yA) -- (\xEnd,\yA);
\draw[stlane] (0.25,\yB) -- (\xEnd,\yB);
\draw[stlane] (0.25,\yC) -- (\xEnd,\yC);
\draw[stlane] (0.25,\yD) -- (\xEnd,\yD);

\node[stsection] at (1.60,3.70) {\textsf{1}};
\node[streqbox] at (1.20,3.12) {\textsc{Accept}$((1{,}0),\, v_a)$};
\node[streqbox] at (1.20,2.52) {\textsc{Accept}$((1{,}0),\, v_b)$};
\draw[streq] (0.85,\yPa) -- (1.55,\yA);
\draw[streq] (0.85,\yPa) -- (1.70,\yB);
\draw[streq] (1.00,\yPb) -- (1.70,\yC);
\draw[streq] (1.00,\yPb) -- (1.85,\yD);
\draw[strsp] (1.55,\yA) -- (2.65,\yPa);
\draw[strsp] (1.70,\yB) -- (2.65,\yPa);
\draw[strsp] (1.70,\yC) -- (2.65,\yPb);
\draw[strsp] (1.85,\yD) -- (2.65,\yPb);
\node[stwarnbox] at (2.45,1.45) {split 2/2};

\draw[dotted, gray] (3.10,-0.16) -- (3.10,3.24);

\node[stsection] at (4.55,3.70) {\textsf{2}};
\node[streqbox] at (4.10,3.12) {\textsc{Prepare}$((2{,}1))$};
\draw[streq] (3.75,\yPa) -- (4.40,\yA);
\draw[streq] (3.75,\yPa) -- (4.55,\yB);
\draw[streq] (3.75,\yPa) -- (4.70,\yC);
\draw[streq] (3.75,\yPa) -- (4.85,\yD);
\draw[strsp] (4.40,\yA) -- (5.75,\yPa);
\draw[strsp] (4.55,\yB) -- (5.75,\yPa);
\draw[strsp] (4.70,\yC) -- (5.75,\yPa);
\node[stnotebox] at (5.45,2.52) {$v_a \times 2,\; v_b \times 1$};

\draw[dotted, gray] (6.15,-0.16) -- (6.15,3.24);

\node[stsection] at (7.45,3.70) {\textsf{3}};
\node[streqbox] at (7.00,3.12) {%
    \textsc{Accept}$((2{,}1),\, v_a)$\\[-1pt]
    \textcolor{blue}{prepare $(3{,}0)$}%
};
\draw[streq] (6.65,\yPa) -- (7.30,\yA);
\draw[streq] (6.65,\yPa) -- (7.45,\yB);
\draw[streq] (6.65,\yPa) -- (7.60,\yC);
\draw[streq] (6.65,\yPa) -- (7.75,\yD);
\draw[strsp] (7.30,\yA) -- (8.55,\yPa);
\draw[strsp] (7.60,\yC) -- (8.55,\yPa);
\node[stcommitbox] at (8.48,\yPa-0.16) {commit $v_a$};

\draw[staxis] (0.25,-0.42) -- node[below, font=\scriptsize] {time} (\xEnd,-0.42);
\end{tikzpicture}%
}% end resizebox
\caption{Fast-round collision and classic recovery. $P_1$ and $P_2$ split the acceptors 2/2 under fast ballot~$(1,0)$, so no value reaches $q_f = 3$. $P_1$ then prepares classic ballot~$(2,1)$, recovers $v_a$ as the unique maximum in the prepare responses, commits it, and piggybacks $(3,0)$ to restore a fast next round.}
\label{fig:collision-recovery}
\end{figure}

% ------------------------------------------------------------------------------
\subsection{Learning Decided Values}
\label{sec:learning-committed-values}
% ------------------------------------------------------------------------------

A proposer may send \textsc{Accept} under ballot $b$ only if it knows the register value for $b$.
It can learn that value by running \textsc{Prepare}, by committing and reusing the piggybacked next ballot, or via a best-effort learn notification from another proposer.
Learn traffic is only an optimization: if it is stale or missing, the proposer falls back to \textsc{Prepare}.
Appendix~\ref{sec:python-implementation} sketches one concrete learn-message format.

% ------------------------------------------------------------------------------

% ==============================================================================
\section{Correctness}
\label{sec:correctness}
% ==============================================================================

Fast CASPaxos must preserve \emph{linearizability}: every completed operation appears to take effect atomically between invocation and response.

\paragraph{Property (First-write-wins).}
An acceptor accepts at most one distinct value per ballot: it accepts the first valid \textsc{Accept} for a fast ballot and rejects any later request carrying a different value at the same ballot.

\paragraph{Lemma 1 (Single decision per ballot).}
At most one value can be decided in any ballot~$b$.
For a classic ballot, only one proposer may use~$b$, so at most one value is proposed.
For a fast ballot, first-write-wins ensures each acceptor accepts at most one value at~$b$; since any two fast quorums intersect, at most one value can reach a fast quorum.

\paragraph{Lemma 2 (Recovery).}
When a proposer runs \textsc{Prepare} at ballot~$b$, it must choose a value consistent with any that could have been decided at a lower ballot.
Any value chosen in a fast ballot was accepted by a fast quorum of $q_f$ acceptors; a prepare quorum of $q_p$ acceptors intersects that fast quorum in at least $q_f + q_p - N \geq \lceil N/4 \rceil + 1$ acceptors, so any chosen value appears in the prepare responses.
A uniquely chosen value therefore appears strictly more often than any competitor; if several values tie for the maximum, none could have been chosen exclusively, and any tied maximum is safe.
In the fixed-quorum setting, this reduces to recovering the unique maximum, or any tied maximum, at the highest fast ballot seen by the prepare quorum~\cite{howard_tr19,howard_arxiv19}.

\paragraph{Theorem (Safety).}
By induction on increasing ballot order: before proposing under ballot~$b$, a proposer must know the register value for~$b$---via \textsc{Prepare}, a reused piggybacked promise, or a learn message.
If no earlier ballot was accepted, any value is safe.
Otherwise, Lemma~1 ensures at most one value was decided at the highest earlier ballot, and Lemma~2 ensures the proposer recovers it.
The proposer applies its update to that value and sends \textsc{Accept} under a higher ballot, so committed register states form a single chain under the update relation, giving linearizability.

\paragraph{Piggybacked promises and duplicate fast accepts.}
A piggybacked prepare only advances promises for later rounds and does not alter the current accept.
Acknowledging a duplicate \textsc{Accept}($b$, $v$) is safe because it introduces no new value at~$b$.

\paragraph{Liveness.}
As with Paxos, FLP precludes guaranteed progress in a fully asynchronous system~\cite{fischer_jacm85}.
In practice, progress requires eventual message delivery and either a stable proposer completing a classic round or a fast round in which only one value is proposed.

\paragraph{TLA+ specification.}
We model Fast CASPaxos in TLA+~\cite{lamport_tla} and check safety and liveness with TLC.
Appendix~\ref{sec:tla-specification} summarizes the model and checked configurations.

% ==============================================================================
\section{Related Work}
\label{sec:related-work}
% ==============================================================================

Fast CASPaxos sits between CASPaxos~\cite{rystsov_caspaxos} and Fast Paxos~\cite{lamport_msr05}: CASPaxos contributes rewritable-register semantics over single-decree Paxos, and Fast Paxos contributes shared fast ballots.
Fast CASPaxos combines them by letting CASPaxos's piggybacked prepare promise either a proposer-owned classic ballot or a shared fast ballot.

The \emph{epochs-by-recovery} framework~\cite{howard_tr19,howard_arxiv19} provides the general theory: multiple proposers may share a ballot provided fast quorums intersect and a strengthened cross-phase intersection lets recovery determine the possible decided values.
Fast CASPaxos is an instance of this framework applied to a rewritable register.
Flexible Paxos~\cite{howard_opodis16} and Fast Flexible Paxos~\cite{howard_icdcn21} show that quorum sizes are a separate design dimension; the recovery rule in Section~\ref{sec:value-recovery} is the natural starting point for a flexible-quorum variant.

EPaxos~\cite{moraru_sosp13} also targets leaderless operation, but for replicated logs rather than a single rewritable register.
Its fast path depends on command commutativity and dependency tracking, whereas Fast CASPaxos treats all writes to one register as conflicting and recovers with ordinary prepare/accept structure.

Rapid~\cite{suresh_atc18} is closer in spirit: it uses Fast Paxos for cluster membership updates in an almost-everywhere-agreement workload where many proposers independently converge on the same value.
Fast CASPaxos benefits from the same pattern via idempotent duplicate fast accepts, but for general read--modify--write register operations.

% ==============================================================================
\section{Summary}
\label{sec:summary}
% ==============================================================================

Fast CASPaxos is a small modification to CASPaxos that yields a large increase in flexibility.
By allowing the piggybacked ballot to be either a shared fast ballot or a proposer-owned classic ballot, the protocol supports both leadered and leaderless operation, and can switch between them on a per-operation basis.

The modification requires only the addition of the Fast Paxos vote-tallying rule for value recovery and the reservation of proposer~ID~0 for shared fast ballots.
Acceptor state remains unchanged---two fields per register---and the protocol retains CASPaxos's simplicity and minimal metadata footprint.

We have formalized the protocol in TLA+ and verified safety and liveness properties through model checking.
Our deterministic simulation evaluation demonstrates that Fast CASPaxos achieves 1~RTT commits in both leadered and leaderless modes when uncontended, and gracefully falls back to coordinated recovery under contention.

More broadly, we believe that adaptive protocols that can tune their behavior to match current conditions---rather than committing to a single operating mode at deployment time---represent a promising direction for practical consensus systems.


% ======================================================================
% 7. BIBLIOGRAPHY
% ======================================================================
\ifIEEE
    \bibliographystyle{IEEEtran}
\fi

\ifACM
    \bibliographystyle{ACM-Reference-Format}
\fi

\ifUSENIX
    \bibliographystyle{plain}
\fi

\bibliography{ref}

% ======================================================================
% APPENDIX
% ======================================================================
\appendix
% ==============================================================================
\section{Formal Specification}
\label{sec:tla-specification}
% ==============================================================================

We formalize Fast CASPaxos in TLA+~\cite{lamport_tla} and check safety and liveness with TLC.
The model includes fast and classic ballots, piggybacked prepare, and fast-ballot recovery by possible decided values.

\paragraph{Model.}
Each acceptor stores \texttt{prepared}, \texttt{accepted}, and \texttt{value}.
Ballots are records $[\text{round} \mapsto r, \text{proposer} \mapsto p]$, with $\text{proposer} = 0$ denoting fast ballots.
To keep the state space tractable, message channels are bounded slots rather than growing sets, and fast accepts are modeled directly in the acceptor action rather than as explicit request messages.

\paragraph{Actions.}
\texttt{PrepareStep} updates a promise and returns accepted state.
\texttt{AcceptReqStep} installs a classic accept request after either piggybacked prepare or classic recovery.
\texttt{AcceptRspStep} processes fast or classic accepts and records the piggybacked prepare.

\paragraph{Checked properties.}
\texttt{Safety} checks typing, single accepted values per classic ballot, single chosen values per ballot, \texttt{accepted} never exceeding \texttt{prepared}, and comparability/reachability of committed values.
\texttt{Transition\-Safety} adds per-transition monotonicity and request/response consistency.
\texttt{Progress} asserts that some non-initial value is eventually committed, using a terminating wrapper to keep liveness checking finite.

\paragraph{Configurations.}
We check safety with 5~acceptors, 3~proposers, 5~values, and 4~rounds, using symmetry reduction and a view projection.
We check liveness with 4~acceptors, 2~proposers, 3~values, and 3~rounds, using weak fairness and a message-count bound.
Both runs report no violations.
The TLA+ source, configurations, and scripts are included in the accompanying repository.

% ==============================================================================
\section{Python Reference Implementation}
\label{sec:python-implementation}
% ==============================================================================

We include a simplified Python implementation of the Fast CASPaxos acceptor and proposer as a pedagogical complement to Figure~\ref{fig:fast-caspaxos-summary}.
A fuller demo harness is available in the accompanying repository.

\paragraph{Ballots and quorums.}
A ballot is a \texttt{(round, proposer)} pair ordered lexicographically; \texttt{proposer\,=\,0} denotes a fast ballot.
Classic quorums require $\lfloor N/2 \rfloor + 1$ acceptors; fast quorums require $\lceil 3N/4 \rceil$.

\paragraph{Acceptor.}
The acceptor handles \textsc{Prepare} and \textsc{Accept} RPCs.
Its state matches Figure~\ref{fig:fast-caspaxos-summary}: \texttt{promised}, \texttt{accepted\_ballot}, and \texttt{accepted\_value}.
The \textsc{Accept} handler uses first-write-wins for fast ballots and processes the piggybacked prepare atomically.

\begin{lstlisting}[
  style=pseudocode,
  numbers=left,
  numberstyle=\tiny\color{gray},
  xleftmargin=2em,
  basicstyle=\scriptsize\ttfamily,
  aboveskip=6pt,
  belowskip=6pt,
]
class Acceptor:
    def __init__(self, name, initial_value=None):
        self.promised = Ballot(1, 0)  # initial fast ballot
        self.accepted_ballot = Ballot(0, 0)
        self.accepted_value = initial_value

    def prepare(self, ballot):                    # Prepare RPC
        if ballot < max(self.promised, self.accepted_ballot):
            return Reject(self.promised)
        self.promised = ballot
        return Promise(self.accepted_ballot, self.accepted_value)

    def accept(self, ballot, value, next_ballot=None):
        # Idempotent replay (same ballot and value)
        if ballot == self.accepted_ballot \
                and value == self.accepted_value:
            return Accepted(self.promised)
        if ballot < max(self.promised, self.accepted_ballot):
            return Reject(self.promised)
        # Same ballot, different value: reject
        if ballot == self.accepted_ballot:
            return Reject(self.promised)
        self.accepted_ballot = ballot
        self.accepted_value = value
        # Piggybacked prepare: prepare next round
        if next_ballot and next_ballot > self.promised:
            self.promised = next_ballot
        return Accepted(self.promised)
\end{lstlisting}

\paragraph{Proposer.}
The proposer runs the Read--Modify--Write cycle.
It reuses a piggybacked prepare when available; otherwise it runs a full \textsc{Prepare}.
The update function is applied locally, and the result is committed with \textsc{Accept}.
With full responses, fast-ballot recovery in \texttt{\_recover} reduces to choosing the unique maximum, or any tied maximum, at the highest fast ballot observed.

\paragraph{Learning.}
The reference implementation models inter-proposer learning as a best-effort notification carrying the decided ballot, value, and prepared next ballot.
If that notification is missing, the proposer falls back to \textsc{Prepare} on its next operation.

\paragraph{Simplifications.}
This reference implementation assumes that all acceptors respond; a production implementation must handle partial responses and retry with a higher classic ballot when fast-ballot recovery remains ambiguous.

\begin{lstlisting}[
  style=pseudocode,
  numbers=left,
  numberstyle=\tiny\color{gray},
  xleftmargin=2em,
  basicstyle=\scriptsize\ttfamily,
  aboveskip=6pt,
  belowskip=6pt,
]
class Proposer:
    def __init__(self, id, acceptors, enable_fast=True):
        self.id = id
        self.acceptors = acceptors
        self.enable_fast = enable_fast
        self.ballot = Ballot(0, 0)
        self.cached_value = None
        self.prepared = None  # piggybacked prepare ballot

    def propose(self, update_fn):
        while True:
            # 1. Read: obtain ballot and current value
            b, need_read = self._choose_ballot()
            if need_read:
                ok, val = self._run_read(b)
                if not ok: continue
                self.cached_value = val
            # 2. Modify: apply the client's update function
            new_val = update_fn(self.cached_value)
            # 3. Write: commit the new value
            if self._run_write(b, new_val):
                self.cached_value = new_val
                return new_val

    def _choose_ballot(self):
        if self.prepared and self.prepared > self.ballot:
            b, self.prepared = self.prepared, None
            self.ballot = b
            return b, False         # reuse piggybacked prepare
        if self.enable_fast \
                and self.ballot < Ballot(1, 0):  # initial fast
            self.ballot = Ballot(1, 0)
            return self.ballot, False  # initial fast ballot
        self.ballot = Ballot(self.ballot.round + 1, self.id)
        return self.ballot, True    # classic: read required

    def _run_read(self, ballot):
        resps = [a.prepare(ballot) for a in self.acceptors]
        promises = [r for r in resps if isinstance(r, Promise)]
        N = len(self.acceptors)
        if len(promises) < N // 2 + 1:
            return False, None
        return True, self._recover(promises, N)

    def _recover(self, promises, N):
        best = max(promises, key=lambda p: p.accepted_ballot)
        if best.accepted_ballot == Ballot(0, 0):
            return self.cached_value
        if best.accepted_ballot.proposer != 0:  # classic
            return best.accepted_value
        # Fast ballot: choose any tied maximum at best ballot
        at_best = [p.accepted_value for p in promises
                   if p.accepted_ballot == best.accepted_ballot]
        counts = Counter(at_best)
        max_count = max(counts.values())
        tied = [v for v, c in counts.items() if c == max_count]
        return tied[0]

    def _run_write(self, ballot, value):
        N = len(self.acceptors)
        next_b = Ballot(ballot.round + 1,
                        0 if self.enable_fast else self.id)
        resps = [a.accept(ballot, value, next_b)
                 for a in self.acceptors]
        accepts = [r for r in resps if isinstance(r, Accepted)]
        needed = -(-3*N//4) if ballot.proposer == 0 \
                 else N//2 + 1
        if len(accepts) >= needed:
            self.prepared = next_b  # cache piggybacked prepare
            return True
        return False
\end{lstlisting}


\end{document}
